\documentclass[border=2pt]{standalone}
\usepackage{polymer-paper-preamble}

\begin{document}
\begin{tikzpicture}[
  x=1cm, y=1cm,
  every node/.style={font=\sffamily\fontsize{7.2}{8.6}\selectfont, text=cGray!95!black},
  panelTitle/.style={font=\sffamily\bfseries\fontsize{8.0}{9.6}\selectfont, text=cTeal, align=center},
  motif/.style={font=\sffamily\fontsize{7.0}{8.4}\selectfont, text=cAmber!90!black, align=center},
  origin/.style={font=\sffamily\fontsize{7.2}{8.6}\selectfont, text=cBrown, align=center},
  originSub/.style={font=\sffamily\fontsize{5.8}{7.0}\selectfont, text=cBrown!85!black, align=center},
  band/.style={font=\sffamily\fontsize{7.2}{8.6}\selectfont, text=cBlue, align=center},
  tickLabel/.style={font=\sffamily\fontsize{4.3}{5.2}\selectfont, text=cGray!90!black},
  evidenceArrow/.style={-{Stealth[length=4pt,width=2.8pt]}, cGray!70, line width=0.8pt},
  sep/.style={cGray!50, line width=0.32pt},
]
\useasboundingbox (0, 0) rectangle (14.0, 7.68);

% ================================================================
% PANEL ELLIPSES
% ================================================================

% Panel 1 — blue ellipse
\draw[cBlue, line width=1.8pt]
  (3.00, 6.28) ellipse [x radius=1.82cm, y radius=1.38cm];

% Panel 2 — orange/amber ellipse
\draw[cAmber!80!black, line width=1.8pt]
  (3.08, 3.98) ellipse [x radius=1.90cm, y radius=0.87cm];

% Panel 3 — purple ellipse
\draw[cBlue!45!cRed, line width=1.8pt]
  (2.84, 1.65) ellipse [x radius=1.78cm, y radius=1.39cm];

% ================================================================
% PANEL 1 — Leakage current evidence (log-log power-law)
% ================================================================
% plot box: x=1.92-3.98, y=5.55-7.19

% Axes
\draw[cGray, line width=0.38pt, -{Stealth[length=2.5pt,width=1.8pt]}]
  (1.92, 5.55) -- (1.92, 7.15);
\draw[cGray, line width=0.38pt, -{Stealth[length=2.5pt,width=1.8pt]}]
  (1.92, 5.55) -- (3.98, 5.55);

% X-axis ticks (7 ticks: 10^-3 ... 10^3)
\foreach \xi/\lab in {
  1.92/{10^{-3}}, 2.263/{10^{-2}}, 2.607/{10^{-1}},
  2.950/{10^{0}}, 3.293/{10^{1}}, 3.637/{10^{2}}, 3.980/{10^{3}}} {
  \draw[cGray, line width=0.25pt] (\xi, 5.55) -- (\xi, 5.49);
  \node[tickLabel, anchor=north] at (\xi, 5.43) {$\lab$};
}

% Y-axis ticks (6 ticks: 10^-4 ... 10^1)
\foreach \yi/\lab in {
  5.55/{10^{-4}}, 5.858/{10^{-3}}, 6.166/{10^{-2}},
  6.474/{10^{-1}}, 6.782/{10^{0}}, 7.090/{10^{1}}} {
  \draw[cGray, line width=0.25pt] (1.92, \yi) -- (1.86, \yi);
  \node[tickLabel, anchor=east] at (1.80, \yi) {$\lab$};
}

% Axis labels
\node[font=\sffamily\fontsize{5.5}{6.6}\selectfont, text=cGray]
  at (2.95, 5.28) {$\log\,t$};
\node[font=\sffamily\fontsize{5.5}{6.6}\selectfont, text=cGray, rotate=90]
  at (1.60, 6.32) {$\log\,I$};

% Power-law straight line (log-log)
\draw[cBlue, line width=0.92pt] (1.92, 7.09) -- (3.98, 5.62);

% Slope annotation (dashed triangle)
\draw[cBlue, dashed, line width=0.45pt] (2.60, 6.48) -- (2.60, 6.10) -- (3.00, 6.10);
\node[font=\sffamily\fontsize{5.8}{7}\selectfont, text=cBlue]
  at (2.45, 6.25) {slope $= -n$};

% I(t) label
\node[font=\sffamily\fontsize{6.5}{7.8}\selectfont, text=cBlue, anchor=west]
  at (2.80, 6.88) {$I(t)\!\propto\!t^{-n}$};

% n arrow (out of panel to right)
\draw[cBlue, line width=0.8pt, -{Stealth[length=3.5pt,width=2.5pt]}]
  (3.60, 6.55) -- (4.60, 6.55);
\node[font=\sffamily\fontsize{7}{8.4}\selectfont, text=cBlue]
  at (4.78, 6.55) {$n$};

% Panel title (drift-gate anchor)
\node[panelTitle] at (0.58, 6.29) {Leakage\\current\\evidence};

% ================================================================
% PANEL 2 — ISPD inference (g(Et) bell curves)
% ================================================================

% Et axis
\draw[cGray, line width=0.38pt, -{Stealth[length=2.5pt,width=1.8pt]}]
  (1.65, 3.10) -- (4.60, 3.10);
\node[font=\sffamily\fontsize{6}{7.2}\selectfont, text=cGray]
  at (4.72, 3.10) {$E_t$};
\node[font=\sffamily\fontsize{6.5}{7.8}\selectfont, text=cGray, rotate=90]
  at (1.45, 3.85) {$g(E_t)$};

% Shallow peak (teal)
\fill[cTeal!25]
  (1.65, 3.12)
  .. controls (1.82,3.15) and (1.95,3.60) .. (2.12, 4.20)
  .. controls (2.29,3.60) and (2.42,3.15) .. (2.62, 3.12)
  -- (1.65, 3.12) -- cycle;
\draw[cTeal, line width=0.9pt]
  (1.65, 3.12)
  .. controls (1.82,3.15) and (1.95,3.60) .. (2.12, 4.20)
  .. controls (2.29,3.60) and (2.42,3.15) .. (2.62, 3.12);

% Deep peak (purple)
\fill[cBlue!45!cRed!22]
  (2.58, 3.12)
  .. controls (2.72,3.15) and (2.85,3.50) .. (3.00, 3.88)
  .. controls (3.15,3.50) and (3.28,3.15) .. (3.48, 3.12)
  -- (2.58, 3.12) -- cycle;
\draw[cBlue!45!cRed, line width=0.9pt]
  (2.58, 3.12)
  .. controls (2.72,3.15) and (2.85,3.50) .. (3.00, 3.88)
  .. controls (3.15,3.50) and (3.28,3.15) .. (3.48, 3.12);

% g(Et) arrow out
\draw[cAmber!80!black, line width=0.8pt, -{Stealth[length=3.5pt,width=2.5pt]}]
  (3.55, 3.75) -- (4.60, 3.75);
\node[font=\sffamily\fontsize{7}{8.4}\selectfont, text=cGray]
  at (4.88, 3.75) {$g(E_t)$};

% Panel title and band labels (drift-gate anchors)
\node[panelTitle] at (0.55, 3.85) {ISPD\\inference};
\node[band] at (1.90, 3.47) {shallow};
\node[band] at (2.83, 3.47) {deep};

% ================================================================
% PANEL 3 — Discharge evidence (Debye decay)
% ================================================================
% plot box: x=1.85-3.83, y=0.97-2.62

% Axes
\draw[cGray, line width=0.38pt, -{Stealth[length=2.5pt,width=1.8pt]}]
  (1.85, 0.97) -- (1.85, 2.55);
\draw[cGray, line width=0.38pt, -{Stealth[length=2.5pt,width=1.8pt]}]
  (1.85, 0.97) -- (3.83, 0.97);

% X-axis ticks
\foreach \xi/\lab in {
  1.85/{10^{-3}}, 2.18/{10^{-2}}, 2.51/{10^{-1}},
  2.84/{10^{0}}, 3.17/{10^{1}}, 3.50/{10^{2}}, 3.83/{10^{3}}} {
  \draw[cGray, line width=0.25pt] (\xi, 0.97) -- (\xi, 0.91);
  \node[tickLabel, anchor=north] at (\xi, 0.85) {$\lab$};
}

% Y-axis ticks
\foreach \yi/\lab in {
  0.97/{10^{-3}}, 1.30/{10^{-2}}, 1.63/{10^{-1}},
  1.96/{10^{0}}, 2.30/{10^{1}}} {
  \draw[cGray, line width=0.25pt] (1.85, \yi) -- (1.79, \yi);
  \node[tickLabel, anchor=east] at (1.73, \yi) {$\lab$};
}

% Axis labels
\node[font=\sffamily\fontsize{5.5}{6.6}\selectfont, text=cGray]
  at (2.84, 0.67) {$\log\,t$};
\node[font=\sffamily\fontsize{5.5}{6.6}\selectfont, text=cGray, rotate=90]
  at (1.52, 1.80) {$\log$ response};

% Debye decay curve (log-log)
\draw[cGray, line width=0.75pt]
  plot[smooth] coordinates {
    (1.85, 2.48) (2.18, 2.35) (2.51, 2.05) (2.84, 1.65) (3.17, 1.30) (3.50, 1.07) (3.83, 0.97)
  };

% tau_d vertical dashed line
\draw[cGray, dashed, line width=0.38pt] (2.84, 0.97) -- (2.84, 1.65);
\node[font=\sffamily\fontsize{5.5}{6.6}\selectfont, text=cGray, anchor=north]
  at (2.84, 0.87) {$\tau_d$};

% tau_d arrow out of panel
\draw[cBlue!45!cRed, line width=0.8pt, -{Stealth[length=3.5pt,width=2.5pt]}]
  (3.40, 1.45) -- (4.40, 1.45);
\node[font=\sffamily\fontsize{7}{8.4}\selectfont, text=cGray]
  at (4.68, 1.45) {$\tau_d$};

% Panel title and Debye label (drift-gate anchors)
\node[panelTitle] at (0.54, 1.68) {Discharge\\evidence};
\node at (2.40, 1.76) {Debye};
\node[font=\sffamily\fontsize{5.8}{7}\selectfont, text=cGray]
  at (2.40, 1.55) {$\exp(-t/\tau)$};

% ================================================================
% PANEL -> MIDDLE ZONE CONVERGENT ARROWS
% ================================================================

% Panel 1 -> top chain (blue)
\draw[cBlue, line width=1.5pt, -{Stealth[length=5pt,width=3.5pt]}]
  (4.81, 6.28) .. controls (5.20, 6.20) and (5.30, 5.80) .. (5.35, 5.55);
\fill[cBlue] (4.81, 6.28) circle (0.08cm);

% Panel 2 -> middle chain (orange)
\draw[cAmber!80!black, line width=1.5pt, -{Stealth[length=5pt,width=3.5pt]}]
  (4.97, 3.98) -- (5.40, 4.20);
\fill[cAmber!80!black] (4.97, 3.98) circle (0.08cm);

% Panel 3 -> bottom chain (purple)
\draw[cBlue!45!cRed, line width=1.5pt, -{Stealth[length=5pt,width=3.5pt]}]
  (4.62, 1.65) .. controls (5.00, 1.70) and (5.20, 2.80) .. (5.35, 3.20);
\fill[cBlue!45!cRed] (4.62, 1.65) circle (0.08cm);

% ================================================================
% MIDDLE ZONE — Polymer chains (3 rows)
% ================================================================

% Chain row 1 (top, S-rich segments, y~5.45)
\draw[cGray, line width=0.85pt]
  plot[smooth] coordinates {
    (5.30,5.45)(5.60,5.62)(5.90,5.45)(6.20,5.62)(6.50,5.45)
    (6.80,5.62)(7.10,5.45)(7.40,5.62)(7.70,5.45)
  };
% S atoms (amber)
\foreach \sx/\sy in {5.60/5.62, 6.20/5.62, 6.80/5.62, 7.40/5.62} {
  \node[text=cAmber, font=\sffamily\fontsize{6.5}{7.8}\selectfont] at (\sx, \sy+0.28) {S};
  \draw[cAmber, fill=white, line width=0.50pt] (\sx, \sy+0.12) circle (0.042cm);
}

% Chain row 2 (middle, y~4.31)
\draw[cGray, line width=0.85pt]
  plot[smooth] coordinates {
    (5.30,4.31)(5.60,4.48)(5.90,4.31)(6.20,4.48)(6.50,4.31)
    (6.80,4.48)(7.10,4.31)(7.40,4.48)(7.70,4.31)(8.00,4.48)(8.30,4.31)(8.60,4.48)(8.90,4.31)
  };
% S atoms
\foreach \sx/\sy in {5.60/4.48, 6.80/4.48, 8.00/4.48} {
  \node[text=cAmber, font=\sffamily\fontsize{6.5}{7.8}\selectfont] at (\sx, \sy+0.28) {S};
  \draw[cAmber, fill=white, line width=0.50pt] (\sx, \sy+0.12) circle (0.042cm);
}
% Trap sites
\foreach \tx/\ty in {6.20/4.48, 7.40/4.48, 8.60/4.48} {
  \node[circle, fill=cBlue!55, draw=cBlue!80!black, inner sep=0pt, minimum size=2mm,
        line width=0.4pt] at (\tx, \ty+0.10) {};
}

% Chain row 3 (bottom, localized traps, y~3.45)
\draw[cGray, line width=0.85pt]
  plot[smooth] coordinates {
    (5.30,3.45)(5.60,3.62)(5.90,3.45)(6.20,3.62)(6.50,3.45)
    (6.80,3.62)(7.10,3.45)(7.40,3.62)(7.70,3.45)(8.00,3.62)(8.30,3.45)(8.60,3.62)(8.90,3.45)
  };
% S atoms
\foreach \sx/\sy in {5.60/3.62, 6.20/3.62, 7.70/3.62, 8.00/3.62, 8.60/3.62} {
  \node[text=cAmber, font=\sffamily\fontsize{6.5}{7.8}\selectfont] at (\sx, \sy+0.28) {S};
  \draw[cAmber, fill=white, line width=0.50pt] (\sx, \sy+0.12) circle (0.042cm);
}
% Localized trap cluster dashed ellipse
\draw[cGray, dashed, line width=0.42pt]
  (7.55, 3.35) ellipse [x radius=0.70cm, y radius=0.35cm];

% Middle zone -> right band diagram arrow
\draw[cGray!80, line width=1.2pt, -{Stealth[length=5pt,width=3.5pt]}]
  (9.40, 4.00) -- (9.85, 4.00);

% ================================================================
% MIDDLE ZONE LABELS (drift-gate anchors — do not move)
% ================================================================

% Dashed bracket under origin annotations
\draw[cGray!60, dashed, line width=0.4pt]
  (5.50, 2.15) -- (9.50, 2.15);
\draw[cGray!60, dashed, line width=0.4pt] (5.50, 2.15) -- (5.50, 2.35);
\draw[cGray!60, dashed, line width=0.4pt] (9.50, 2.15) -- (9.50, 2.35);
\draw[cGray!60, dashed, line width=0.4pt] (7.50, 2.15) -- (7.50, 1.55);

\node[motif] at (7.35, 6.02) {S-rich segments};
% S-rich segments -> chain S atom arrows
\draw[cGray!70, line width=0.4pt, -{Stealth[length=2.5pt,width=1.8pt]}]
  (6.50, 5.93) -- (6.22, 5.76);
\draw[cGray!70, line width=0.4pt, -{Stealth[length=2.5pt,width=1.8pt]}]
  (7.35, 5.93) -- (7.35, 5.76);
\draw[cGray!70, line width=0.4pt, -{Stealth[length=2.5pt,width=1.8pt]}]
  (8.20, 5.93) -- (7.58, 5.76);
\node[motif] at (7.21, 1.37) {localized traps};
\node[origin] at (6.22, 2.74) {chemical origin};
\node[originSub] at (6.24, 2.40) {(electronegativity, polarizability of S)};
\node[origin] at (8.35, 2.74) {physical origin};
\node[originSub] at (8.36, 2.38) {(local potential fluctuations)};

% ================================================================
% RIGHT ZONE — Band diagram (teal border + CB/VB + g(Et))
% ================================================================

% Teal border
\draw[cTeal, line width=1.2pt, rounded corners=4pt]
  (9.89, 1.89) rectangle (13.78, 6.38);

% CB bar
\fill[cLGray!40] (10.26, 5.69) rectangle (11.52, 5.89);
\draw[cGray!60, line width=0.4pt] (10.26, 5.69) rectangle (11.52, 5.89);
\node[font=\sffamily\bfseries\fontsize{7}{8.4}\selectfont, text=cGray!80, anchor=center]
  at (10.89, 5.79) {CB};

% VB bar
\fill[cLGray!40] (10.27, 2.23) rectangle (11.47, 2.42);
\draw[cGray!60, line width=0.4pt] (10.27, 2.23) rectangle (11.47, 2.42);
\node[font=\sffamily\bfseries\fontsize{7}{8.4}\selectfont, text=cGray!80, anchor=center]
  at (10.87, 2.32) {VB};

% Energy axis
\draw[cGray, line width=0.7pt, -{Stealth[length=3.2pt,width=2.4pt]}]
  (11.20, 2.42) -- (11.20, 5.69);
\node[font=\sffamily\fontsize{6.5}{7.8}\selectfont, text=cGray, rotate=90, anchor=south]
  at (11.02, 4.05) {Energy};

% Et dashed vertical line
\draw[cGray!45, dashed, line width=0.38pt]
  (12.10, 2.05) -- (12.10, 6.20);
\node[font=\sffamily\fontsize{7}{8.4}\selectfont, text=cGray]
  at (12.30, 4.15) {$E_t$};

% Shallow trap levels (amber, just below CB)
\foreach \yi in {5.42, 5.22, 5.02, 4.82} {
  \draw[cAmber, line width=0.72pt] (10.55, \yi) -- (10.88, \yi);
  \node[circle, fill=cBlue!50, draw=cBlue!80!black, inner sep=0pt, minimum size=1.4mm,
        line width=0.35pt] at (10.71, \yi+0.07) {};
}
\node[text=cAmber] at (11.08, 5.12) {$\cdots$};

% Deep trap levels (purple, just above VB)
\foreach \yi in {3.20, 2.95} {
  \draw[cBlue!45!cRed, line width=0.72pt] (10.55, \yi) -- (10.88, \yi);
  \node[circle, fill=cBlue!50, draw=cBlue!80!black, inner sep=0pt, minimum size=1.4mm,
        line width=0.35pt] at (10.71, \yi+0.07) {};
}
\node[text=cBlue!45!cRed] at (11.08, 3.05) {$\cdots$};

% g(Et) vertical axis and horizontal axis
\draw[cGray, line width=0.6pt, -{Stealth[length=2.5pt,width=1.8pt]}]
  (12.25, 2.05) -- (12.25, 6.00);
\draw[cGray, line width=0.6pt, -{Stealth[length=2.5pt,width=1.8pt]}]
  (12.25, 2.05) -- (13.70, 2.05);
\node[font=\sffamily\fontsize{7}{8.4}\selectfont, text=cGray, anchor=south]
  at (12.98, 6.05) {$g(E_t)$};
\node[font=\sffamily\fontsize{6}{7.2}\selectfont, text=cGray]
  at (13.10, 1.75) {$g(E_t)$};

% Shallow g(Et) Gaussian (peach/amber, x=12.25-12.95, y=4.17-5.67)
\fill[cAmber!28]
  (12.25, 4.17)
  .. controls (12.52, 4.25) and (12.75, 4.70) .. (12.75, 4.92)
  .. controls (12.75, 5.15) and (12.52, 5.58) .. (12.25, 5.67)
  -- cycle;
\draw[cAmber!80!black, line width=0.8pt]
  (12.25, 4.17)
  .. controls (12.52, 4.25) and (12.75, 4.70) .. (12.75, 4.92)
  .. controls (12.75, 5.15) and (12.52, 5.58) .. (12.25, 5.67);
% shallow label omitted here — pdftotext extracts this before the ISPD-panel
% "shallow" node, causing a drift-gate artifact. Color (cAmber) identifies it visually.

% Deep g(Et) Gaussian (purple, x=12.25-13.03, y=2.91-3.98)
\fill[cBlue!45!cRed!22]
  (12.25, 2.91)
  .. controls (12.55, 2.99) and (12.80, 3.28) .. (12.80, 3.44)
  .. controls (12.80, 3.62) and (12.55, 3.90) .. (12.25, 3.98)
  -- cycle;
\draw[cBlue!45!cRed, line width=0.8pt]
  (12.25, 2.91)
  .. controls (12.55, 2.99) and (12.80, 3.28) .. (12.80, 3.44)
  .. controls (12.80, 3.62) and (12.55, 3.90) .. (12.25, 3.98);
% deep label omitted here — same reason as shallow above.

% Converged trap-depth picture arrow + label
\draw[cTeal, line width=0.9pt, -{Stealth[length=4pt,width=2.5pt]}]
  (11.85, 1.89) -- (11.85, 1.40);
\node[font=\sffamily\fontsize{7}{8.4}\selectfont, text=cTeal,
      anchor=north, align=center]
  at (11.85, 1.32) {converged\\trap-depth picture};

\end{tikzpicture}
\end{document}
